% =====================================================================
%  Section: Confinement
%  Paper 1, "It from Bit via Gödel". Drafted 2026-06-13 from the plan
%  confinement_section_plan_2026-06-13.md, with the unification decided
%  in discussion 2026-06-13: ALL mass is associator-geodesic length.
%
%  DEFINES sec:confinement. Sited after the GHZ->colour / W->lepton
%  proofs and the generations material, which it uses.
%
%  CONVENTIONS: British English; \autocite (biblatex); amsthm with the
%  paper-wide openproblem environment (thmtools, declared in preamble);
%  status-flagged prose; NO bullet lists in prose. The K3 four-class
%  enumeration is a TABLE (notebook I retired -- it was pure combinatorics,
%  no quantum circuit; the enumeration lives here as a table, with a tiny
%  scripts/ .py for reproducibility in the repo).
%
%  THREE STANDALONE NUMBERED OPEN PROBLEMS (per JS 2026-06-13):
%    op:confinement-potential   -- V(r) as "the associator-geodesic length
%                                  resolved by emergent distance" (links to
%                                  op:associator-geodesic structurally;
%                                  emergent-distance part = Paper 2, NOT a
%                                  blocker for Paper 1)
%    op:confinement-energy-locus -- quark triple vs gluon Bell pairs;
%                                  "no medium" committed, gluon-Bell-pair
%                                  location not excluded; static, attackable
%                                  with existing entanglement tools
%    op:glueball-spectrum       -- the four-class ggg spectrum + 0++ mass;
%                                  mass origin = associator-geodesic (all
%                                  mass is), on C-only closed networks
%
%  UNIFICATION (committed 2026-06-13): all mass -- fermion, glueball,
%  baryon binding -- is associator-geodesic length. Baryon mass is mostly
%  binding (interaction) energy, so interaction energy IS associator-
%  geodesic length, which is why the confinement potential and the mass
%  map are the SAME object resolved two ways. op:associator-geodesic is
%  hereby the mass map for ALL mass, not just fermions.
%
%  Single forward sentence on gravity + E=mc^2 (#4,#5 from discussion),
%  full discussion deferred to Paper 2.
%
%  CROSS-REFS OUT (verify all before compile):
%    sec:entanglement-classes, sec:generations, sec:register,
%    op:associator-geodesic, sec:predictions, app:qiskit-composition,
%    app:qiskit-meson, app:qiskit-handles.
%    sec:space-of-computations + sec:gr-nonlocality (RESOLVED 2026-06-13:
%    the spacetime section's label is sec:space-of-computations; the
%    pre-spatial subsection is sec:gr-nonlocality).
%  NEW CITATIONS (RIS to prepare): coffmankunduwootters2000 (three-tangle),
%    osborneverstraete2006 (general monogamy). Check if already in .bib.
%  NOTEBOOKS quoted: G (notebook_G_confinement_budget),
%    H (notebook_H_glueball_fanout). Notebook I retired.
% =====================================================================

\section{Confinement}
\label{sec:confinement}

Confinement is the statement that colour-charged particles are never found in isolation: no free
quark, no free gluon, only colour-singlet combinations. In the standard account this is attributed
to the dynamics of the non-Abelian colour field --- the flux tube, the linear potential, the running
coupling --- and remains, in its quantitative form, among the hardest open problems in field theory.
The framework offers a different account of \emph{why} confinement holds, one that is structural and
exact rather than dynamical and asymptotic, and it is careful about the boundary between what that
account settles and what it leaves open. We place the section here, beside the GHZ$\to$colour and
W$\to$lepton results (\S\ref{sec:entanglement-classes}, \S\ref{sec:generations}) it depends on,
because confinement in this framework is the same kind of fact as those: a consequence of the
entanglement class, read off the structure rather than computed from a force.

\subsection{Confinement as an exhausted entanglement budget}

A baryon's three quarks are bound through their colour ($C$) qubits, and the binding state is the
maximal GHZ state --- the colour singlet. The decisive property of that state is supplied by the
monogamy of entanglement. Entanglement cannot be freely shared: a system that is maximally entangled
with some partners has none left to spend on others. Quantitatively, for three qubits the
Coffman--Kundu--Wootters monogamy inequality
\autocite{coffmankunduwootters2000distributed,OsborneVerstraete2006} bounds the total,
\begin{equation}
  \tau_{A(BC)} \;\ge\; \tau_{AB} + \tau_{AC},
\end{equation}
and the genuinely tripartite part --- the three-tangle $\tau_3$ --- measures what is held
irreducibly among all three. For the maximal GHZ state $\tau_3 = 1$: the budget is fully spent on
the tripartite binding, and \emph{every} bipartite tangle, internal or external, is forced to zero.

The consequence is confinement. A colour singlet has $\tau_3 = 1$, so it has no entanglement
whatever left over to share with anything outside the triple. A quark cannot be entangled with an
external colour charge because the singlet it belongs to has spent its entire colour-entanglement
budget on its two partners. There is no free quark because there is no budget for the external
colour correlation a free quark would require. This is verified directly in
Appendix~\ref{app:qiskit-composition} and its companion computation: the colour singlet returns
$\tau_3 = 1$ with both bipartite tangles identically zero, and an attempt to entangle one quark with
an external qubit succeeds only in lockstep with the destruction of the singlet --- the external
correlation rises exactly as the GHZ fidelity falls, so a quark can acquire an outside partner only
by ceasing to be part of a colour singlet.

The reading deserves emphasis because it differs from the standard one in kind, not merely in
technique. This is not a force that grows with separation; it is a property of the \emph{state},
with no distance in it. The three-tangle $\tau_3$ is a function of the amplitudes, and a GHZ state
has $\tau_3 = 1$ whether its qubits are adjacent or galaxies apart --- entanglement is not
distance-dependent, and nothing in the argument above mentions a separation. Confinement, in this
account, is the flat structural fact that a colour singlet has no entanglement available for an
external colour partner. It is also, in consequence, \emph{exact}: a budget is either saturated or
it is not, and the maximal GHZ state saturates it completely. Where lattice QCD extracts confinement
as an asymptotic, numerically-determined property of the colour field, the framework states it as an
exact property of the singlet state.

\subsection{One structure, two consequences: confinement and the vector-like colour force}

The same feature of the GHZ state that gives confinement also explains, at one stroke, why the
colour force is not chiral. The GHZ state has no oriented pairwise substructure: its entanglement is
irreducibly tripartite, carried among all three qubits at once, with --- as the budget shows --- no
bipartite content to host a relative orientation. In \S\ref{sec:generations} this was what forced
the colour sector to be vector-like and the QCD vacuum angle to vanish structurally, there being no
oriented pairwise links to carry a CP-violating phase. It is the same fact, seen from the other side,
that gives confinement: \emph{all} of the colour entanglement is tripartite and internal, which is
at once why none is available externally (confinement) and why none is pairwise and orientable
(chirality-freedom). Confinement and the vector-like character of the strong force are not two
properties of the colour sector that happen to coexist; they are one property of the GHZ class read
twice. The contrast with the W (lepton) sector, whose entanglement \emph{is} pairwise and which is
correspondingly chiral and unconfined, is the same contrast in reverse.

\subsection{The confinement potential is distance-free in origin}

The structural account establishes \emph{that} quarks cannot be isolated; it does not yet yield the
\emph{measured, distance-dependent} phenomenology --- the static quark--antiquark potential, its
linear rise, the slope that lattice QCD and Regge trajectories call the string tension. We state the
gap precisely, because it is more interesting than a missing number. Distance, in this framework, is
not fundamental: it is emergent and relational, built from the network of interactions
(\S\ref{sec:space-of-computations}). So the task is not to compute the energy of a flux tube stretched across a
pre-existing space; it is to show how a potential that \emph{depends on separation} arises at all
from a constraint that contains no separation.

The unification of the mass sector (below) fixes the form this answer must take. If, as the framework
holds, all mass is associator-geodesic length, then the interaction energy binding quarks --- which
is most of a hadron's mass --- is associator-geodesic length, and the confinement potential is that
length resolved by separation. The confinement potential and the mass map are therefore not
independent problems: the potential is what the geodesic length looks like when one asks how it
varies with the emergent distance between the charges.

\begin{openproblem}[The confinement potential]\label{op:confinement-potential}
Derive the static quark--antiquark potential --- equivalently the quantity lattice QCD and Regge
trajectories call the string tension --- as the associator-geodesic length of
Open Problem~\ref{op:associator-geodesic}, resolved by the emergent distance of
\S\ref{sec:space-of-computations}, and show that it reproduces asymptotic freedom at short range as the limit in
which little tripartite budget is yet committed. No flux tube or confining medium is posited; the
content is that the confining \emph{energy} is the same geodesic length that fixes mass, expressed as
a function of relational separation. The emergent-distance half of this engages the spacetime
programme of Paper~2 and is not required for the structural confinement result above.
\end{openproblem}

There is a further question the structural account raises but does not settle: \emph{where} the
confining energy resides. The monogamy obstruction shows that a colour singlet can host no external
correlation, but it does not, by itself, say whether the energy of binding is to be read as a
property of the quark triple's own tripartite entanglement or of the gluon Bell-pair network that
mediates the colour field between them. We commit to the position that there is no confining
\emph{medium} --- no background condensate doing the work, the entanglement being intrinsic and
relational rather than a property of the vacuum --- but that commitment does not decide between the
two intrinsic locations, and we do not exclude that the binding energy is carried by the gluon Bell
pairs rather than the quark triple.

\begin{openproblem}[Where the confining energy resides]\label{op:confinement-energy-locus}
Determine whether the energy of colour binding is carried by the tripartite entanglement of the
quark $C$-qubits or by the gluon Bell-pair network that mediates the colour field, given that it is
in either case intrinsic and relational and not the property of any vacuum medium. Unlike
Open Problem~\ref{op:confinement-potential}, this concerns the static structure of the singlet and
is attackable with the entanglement measures already in hand, without the emergent-distance
machinery.
\end{openproblem}

\subsection{Glueballs}

Glueballs are bound states of the colour field with no valence quarks: in the framework's language,
networks of gluon Bell pairs on $C$-qubits, closed on themselves so that the outer endpoints
contract to a colour singlet with no external fermion register to terminate on. They carry a subtlety
the quark sector does not, and it is the deepest point of the section.

A quark trio produces a \emph{classical} colour record for free, because the GHZ structure is special
Frobenius --- it possesses the copy operation that broadcasts a record (the FANOUT of
\S\ref{sec:entanglement-classes}). A single gluon cannot: its Bell pair is not a copy structure and
has nothing to broadcast. For a glueball to be a real, classically detectable bound state, the gluon
network must therefore \emph{manufacture} the copy structure through its closure topology, asking the
geometry of the network to do what a quark trio does intrinsically. Appendix~\ref{app:qiskit-meson}
and its companion computation make this concrete, and the result is sharper than the question. The
two-gluon network does not, on its open legs, reconstruct the copy structure: it returns a Bell pair.
The closed three-gluon triangle returns exactly $\mathrm{GHZ}_3$, the copy structure itself. The
closure manufactures the record-broadcasting capacity that the constituents lack --- but it takes the
three-gluon triangle to do it, not the two-gluon scalar. The framework's reading is, in consequence,
that classically detectable glueball structure is tied to the three-gluon topology.

The three-gluon glueball carries a structural prediction that goes beyond standard QCD. Three gluons
joined as a triangle form a $K_3$ network --- the same complete graph on three vertices that is not
two-colourable, the obstruction that in \S\ref{sec:generations} generates the three fermion
generations. The four resolution classes of $K_3$ --- three frustrated classes and one symmetric ---
should therefore appear as four channel classes for three-gluon glueballs:

\begin{center}
\begin{tabular}{lll}
\hline
class & orientation assignments & character \\
\hline
frustrated on edge $AB$ & $\{001,\,110\}$ & one edge forced like-oriented \\
frustrated on edge $BC$ & $\{011,\,100\}$ & one edge forced like-oriented \\
frustrated on edge $CA$ & $\{010,\,101\}$ & one edge forced like-oriented \\
symmetric               & $\{000,\,111\}$ & all edges like-oriented \\
\hline
\end{tabular}
\end{center}

Standard QCD organises three-gluon singlets by two couplings, the symmetric $d^{abc}$ and the
antisymmetric $f^{abc}$. The framework predicts four channel classes where QCD sees two --- a
differential, lattice-checkable statement, and the same $K_3$ obstruction that gives the generations,
now in the colour field. That the obstruction reproduces a four-class structure independently in the
gluon sector is itself evidence that it is a genuine organising principle and not a coincidence of
the fermion counting.

\begin{openproblem}[The glueball spectrum]\label{op:glueball-spectrum}
Derive the three-gluon four-class spectrum above and the mass of the lightest ($J^{PC} = 0^{++}$)
glueball as the associator-geodesic length (Open Problem~\ref{op:associator-geodesic}) of the closed
$C$-only Bell network, the lattice value being $\approx 1.7\ \mathrm{GeV}$. The glueball mass shares
its origin with all other mass in the framework --- it is geodesic length --- but on a network living
entirely in the colour sector, with none of the hypercharge or weak-isospin activity that the fermion
masses carry, so the relevant geometry is that of the pure-colour closed network rather than the
mixed-fibre case of the charged fermions.
\end{openproblem}

\subsection{All mass is geodesic length}

The glueball mass is not a special case but an instance of a single principle the framework now
commits to: \emph{all} mass is associator-geodesic length. The masses of the individual quarks and
leptons (\S\ref{sec:generations}), the masses of glueballs, and --- the demanding case --- the masses
of hadrons are all the same geometric quantity, the length of the relevant closed network read
through the octonionic associator. The demanding case is demanding because most of a hadron's mass is
not the mass of its constituents: the proton weighs about $938\ \mathrm{MeV}$, of which the
current-quark masses contribute some nine, the remainder being the energy of the colour binding. Under
the committed principle this binding energy must itself be associator-geodesic length --- the geodesic
length of the closed three-quark colour network --- of which the valence-quark contribution is a per
cent and the rest is the geometry of the binding. This is a sharp and falsifiable burden: the
framework claims the proton mass is the associator-geodesic length of the closed $uud$ network, a
claim that meets the one number lattice QCD has determined most precisely. It also explains why the
framework, which reproduces the charged-lepton masses to better than a per cent, does not yet produce
baryon masses: the obstacle is not the constituent masses but the binding geometry, and the binding is
where almost all of the mass lives. That single geometric quantity --- the associator-geodesic length
of Open Problem~\ref{op:associator-geodesic} --- now carries the entire mass sector, which is a real
localisation of the problem (mass is one thing, not many) with one large computation owed behind it.

One forward-looking consequence deserves a single sentence here and its proper treatment in Paper~2:
because mass is the gravitational charge, the claim that all mass --- including the binding energy that
is most of it --- is associator-geodesic length means that what gravitates is associator-geodesic
length, so that the equivalence of energy and mass and the coupling of mass to spacetime must both, in
the end, be statements about this one geometric quantity.

% =====================================================================
%  SUBSECTION: The Yang--Mills mass gap
%  Drafted 2026-07-03 per JS: having committed to all mass as
%  associator-geodesic length (including the proton's), the paper must
%  say why it does not extend the claim to the Clay Millennium problem
%  -- "it would seem strange to get so far and not make the extra
%  effort" -- with an explanation and a forward-looking suggestion.
%
%  PLACEMENT: into confinement_section_2026-06-13.tex, immediately
%  AFTER the "All mass is geodesic length" subsection (i.e. after its
%  closing gravity/E=mc^2 forward sentence, the paragraph beginning
%  "One forward-looking consequence deserves a single sentence...")
%  and BEFORE \subsection{Confinement and pre-spatial co-location}.
%
%  REGISTER (deliberate): this is the highest-crankery-risk topic in
%  physics; the subsection names the problem precisely so that no claim
%  on it can be read into the section by implication, declines it with
%  specific reasons (same kinematic/dynamical boundary as the QED
%  scope statement), then records the framework-native structural seed
%  (conditional, flagged) and two forward directions. NO new open
%  problem is minted: the framework's pieces of the gap are already
%  op:associator-geodesic, op:glueball-spectrum, op:qcd-scale,
%  op:confinement-potential, and the four translation interfaces are
%  given as a prose programme.
%
%  DEPENDENCY: references op:qcd-scale, which is DEFINED in
%  associator_debt_higgs_mechanism_2026-07-03.tex -- paste that
%  subsection into the boson-masses section first, or this reference
%  dangles. (Fallback if you prefer: drop the op:qcd-scale clause.)
%
%  CROSS-REFS (all existing): sec:interactions, sec:space-of-computations,
%  sec:masses, op:associator-geodesic, op:glueball-spectrum,
%  op:confinement-potential, op:qcd-scale [see dependency].
%  NEW citation keys (RIS in mass_gap_refs_2026-07-03.ris):
%  jaffewitten2006, cubitt2015spectralgap.
%
%  OPTIONAL one-line addition to "What is settled and what is owed",
%  end of the owed paragraph:
%    "The Yang--Mills mass gap is named in \S\ref{sec:mass-gap}
%     precisely so that no claim on it is read into this section by
%     implication."
% =====================================================================

\subsection{The Yang--Mills mass gap}
\label{sec:mass-gap}

A reader who has followed this section to here will have noticed that its commitments walk up to
the edge of a famous problem, and it would be evasive not to name it. The Clay Millennium
formulation asks for two things \autocite{jaffewitten2006}: a mathematically rigorous construction
of quantum Yang--Mills theory on $\mathbb{R}^{4}$, for any compact simple gauge group, meeting the
Wightman axioms or standards of comparable stringency; and a proof that the constructed theory has
a mass gap $\Delta > 0$ --- that its lightest excitation is strictly massive. As physics the gap is
not in doubt: lattice computation exhibits it, and its lightest carrier is the scalar glueball near
$1.7$ GeV whose derivation this section has already recorded as
Open Problem~\ref{op:glueball-spectrum}. The Millennium problem is a problem of proof, not of
phenomenology. This section has asserted that confinement is exact and structural, that glueballs
are closed colour networks, and that all mass --- theirs included --- is associator-geodesic length.
Having gone that far, we owe the reader a plain statement of why we go no further, and of what
going further would require.

The reasons are the boundary already drawn for quantum electrodynamics in \S\ref{sec:interactions},
here made specific. The mass gap is a statement about the spectrum of a quantum dynamics, and this
framework has not constructed that dynamics: there is no Hamiltonian here whose spectrum could be
bounded, no Hilbert space of the field theory, and no continuum $\mathbb{R}^{4}$ to put one on ---
spacetime itself being emergent in this framework, its construction deferred
(\S\ref{sec:space-of-computations}). One cannot prove a spectral gap of an operator one has not
built. Beneath that lies the framework's own open centre: the map from network geometry to physical
mass is Open Problem~\ref{op:associator-geodesic}, and the statement ``every closed colour network
has strictly positive geodesic length'' cannot be a theorem before the functional it concerns
exists. Even were both supplied, the translation into the Millennium problem's terms crosses four
interfaces, each a research programme in its own right: the continuum construction of the emergent
spacetime; the translation of the categorical and entanglement-theoretic statements of this paper
into an axiomatic form; the passage from the algebraic obstruction of this section to the standard
criteria of confinement, the Wilson-loop area law above all; and the spectral proof itself. We claim
none of these. One further limit of scope deserves its own sentence: the Millennium problem is posed
for every compact simple gauge group, while the structural story of this section is specific to the
colour sector as it arises here, on the octonionic register --- our route, even if completed, would
speak to $SU(3)$, not to the problem in its full generality.

What the framework does supply, and what we record because it gives the gap a concrete identity, is
a structural reason to expect one. The photon is massless because the electroweak sector contains a
debt-free direction: the associator-blind $\mathbb{C}$, along which nothing is owed
(\S\ref{sec:masses}). The colour sector contains no such direction --- colour lives on the register
exactly where the associator is non-trivial --- so no closed colour network can be debt-free. And
closed networks are topologically discrete: the minimal closure is a finite structure, and there is
no continuum of ever-smaller closed networks by which geodesic lengths could accumulate towards
zero. If, as the framework's mass formulas everywhere suggest, the geodesic length of a closed
network is fixed by its discrete topology rather than by any continuously variable modulus, then the
spectrum of pure-colour states is bounded below by the length of the minimal closed network. On this
reading the gap exists because the colour sector has no massless direction and no arbitrarily small
closed structure; and $\Delta$ is then not an abstract bound but a particular number --- the mass of
the lightest glueball, precisely the object of Open Problem~\ref{op:glueball-spectrum}, with its
scale in physical units the object of Open Problem~\ref{op:qcd-scale}. Everything in this paragraph
is conditional on the open map: it is a heuristic the framework makes natural, not a theorem the
framework possesses, and we flag it as such.

Two forward directions are worth leaving on the table. The first reframes the problem's notorious
difficulty. Cubitt, P\'{e}rez-Garc\'{i}a and Wolf proved that the spectral-gap question is, in
general, undecidable: for translationally invariant lattice Hamiltonians it encodes the halting
problem, and their own statement of the result names the Yang--Mills gap conjecture as a particular
case of the general problem so shown undecidable \autocite{cubitt2015spectralgap}. For a framework
whose spine is self-referential non-computability this is orientation rather than discouragement.
It says that any proof must exploit special structure --- no general algorithmic route exists ---
and it locates the difficulty exactly where the bump-in-the-carpet principle of this framework
expects to find it. The suggestion we record, in the same register as the other prospects of this
paper --- a direction with a possible path, not a claim --- is that the special structure is the
division-algebra rigidity itself, the Hurwitz--Adams rigidity that terminated the Hopf tower, and
that a proof along this route would show the geodesic functional on closed colour networks to be
bounded below by the minimal-network length, the discreteness of network topology doing the work
that lattice regularisation does elsewhere. The second direction is nearer, and does not wait on the
emergent-distance machinery: in a fixed lattice geometry, the question of whether the
budget-saturation constraint of this section reproduces area-law Wilson loops is one that
tensor-network methods can address now, and it is the interface at which this framework's account of
confinement first becomes commensurable with the standard one. We leave the problem, and the route,
in the record of open problems rather than in the claims of this paper.


\subsection{Confinement and pre-spatial co-location}

A final reading ties confinement to the framework's account of space, and anticipates the fuller
treatment of \S\ref{sec:gr-nonlocality}. In that account entanglement is the pre-spatial relational
structure: entangled systems are not at distinct locations that some influence must bridge, but share
a relation that was never spatially separated to begin with, acquiring distinct locations only by
interacting with other systems that already hold locations in the network. Confinement is exactly
this seen in the colour sector. The three quarks of a singlet are maximally entangled, so in the
pre-spatial structure they are not separated; and the colour-singlet constraint, which forbids them
the external entanglement that location-acquisition requires, is precisely what prevents any one of
them from acquiring a distinct location of its own. A quark cannot be pulled to a place apart because,
in the only structure that is fundamental, it has no place apart to be pulled to, and the singlet will
not let it acquire one. The distance-free monogamy obstruction of this section and the pre-spatial
co-location of \S\ref{sec:gr-nonlocality} are one fact: confinement is the impossibility of giving a
colour charge a separate location while its colour entanglement is wholly spent on its partners.

\subsection{What is settled and what is owed}

Settled, and structural: confinement holds, as the saturation of the tripartite entanglement budget
by the colour singlet; it is exact rather than asymptotic; it is the same fact as the vector-like,
CP-conserving character of the colour force; glueballs are closed gluon networks that must manufacture
their own copy structure through closure, which ties classically detectable glueball states to the
three-gluon topology; and the three-gluon sector carries a four-class structure where standard QCD has
two couplings.

Owed, and recorded as the open problems above: the distance-dependent confinement potential, as the
associator-geodesic length resolved by emergent distance (Open Problem~\ref{op:confinement-potential});
the location of the confining energy, quark triple or gluon network (Open
Problem~\ref{op:confinement-energy-locus}); and the glueball spectrum and lightest mass (Open
Problem~\ref{op:glueball-spectrum}). All three sit downstream of the single mass map, Open
Problem~\ref{op:associator-geodesic}, which the unification of this section makes responsible for the
whole of mass --- the fermion spectrum, the glueball spectrum, and the hadronic binding energy
alike. The Yang--Mills mass gap is named in \S\ref{sec:mass-gap} precisely so that no claim on it is
read into this section by implication.
